\documentclass[conference]{IEEEtran}
\usepackage{cite}
\usepackage{amsmath,amssymb,amsfonts}
\usepackage{array}
\usepackage{graphicx}
\usepackage{textcomp}
\usepackage{xcolor}
\usepackage{balance}

\newcommand{\traceres}{\textsc{TraceRes}}

\begin{document}

\title{Is There a Diffraction Limit?\\
A Task-Specific Evidence Proposal for Credible Optical Nanoscopy}

\author{\IEEEauthorblockN{Research Plan Artifact}
\IEEEauthorblockA{Evidence-bounded Survey, Idea, ExperimentDesign, and Author workflow}}

\maketitle

\begin{abstract}
The phrase ``diffraction limit'' is frequently used as though it were a universal numerical barrier for every optical measurement. That interpretation is too broad. Classical diffraction sets a well-defined resolution baseline for a specified linear far-field imaging model, finite numerical aperture, wavelength, contrast criterion, and measurement objective. Modern super-resolution microscopy changes the inference problem by adding controlled fluorophore-state transitions, patterned illumination, temporal sparsity, targeted excitation, physical expansion, or explicit reconstruction priors. These resources can support information below a conventional Rayleigh separation scale, but they do not erase limits imposed by photon statistics, background, sampling density, labeling linkage, drift, aberration, model mismatch, dynamics, and calibration transfer. This paper proposes \traceres{}, a task-specific resolution passport for converting a nominal nanoscale image into an auditable scientific claim. The proposal separates diffraction-bounded imaging, effective PSF confinement, localization precision, structural resolution, and coordinate accuracy. It specifies blinded synthetic and physical calibration cases, an error-bound envelope, modality-aware counterfactuals, and bounded claim statuses. The artifact is a research design only; no microscope acquisition, reconstruction run, biological measurement, or benchmark result is reported.
\end{abstract}

\begin{IEEEkeywords}
diffraction limit, super-resolution microscopy, STED, SIM, single-molecule localization microscopy, MINFLUX, resolution validation, optical nanoscopy
\end{IEEEkeywords}

\section{Introduction}

The 2014 Nobel Prize in Physics recognized methods that brought optical microscopy beyond the resolution ordinarily associated with classical far-field diffraction. The achievement did not make diffraction irrelevant. Instead, it exposed a more precise scientific question: which quantity is limited, by which physical model, and which additional information resource permits a measurement to distinguish a finer feature? A conventional image may have a diffraction-broadened point-spread function (PSF), while a sparse fluorescent emitter can be localized with much finer statistical precision. A reconstructed image may contain narrow lines, while its labeling density or drift uncertainty prevents the corresponding biological structure from being resolved. A method may demonstrate a small scale on a calibration object yet fail to transfer that claim to a dynamic, scattering, or densely labeled sample.

The literature describes multiple mechanisms for super-resolution optical microscopy: patterned illumination, nonlinear fluorescent-state control, sparse temporal activation, targeted localization, sequential imaging, physical expansion, and computational reconstruction \cite{hell2015,huang2010}. These methods are scientifically powerful because they use more structure than the conventional intensity image. Their diversity also makes a single, unconditional answer to ``is there a diffraction limit?'' unhelpful.

This report develops an evidence-driven proposal called \traceres{}. Its purpose is not to rank instruments by a smallest advertised number. Its purpose is to require a claim to name its task, information channel, error budget, calibration scope, and remaining alternatives. The central contribution is a falsifiable passport that reports what kind of resolution has been validated and what has not. The resulting research question is: \emph{under which definitions and experimentally verifiable conditions can optical microscopy support structural information below the conventional diffraction-limited separation scale?}

\section{What Classical Diffraction Does and Does Not State}

\subsection{A conditional optical baseline}

For two equal incoherent point sources in a conventional far-field system, a commonly used lateral Rayleigh scale is
\begin{equation}
 \delta_R \approx \frac{0.61\lambda}{\mathrm{NA}},
\label{eq:rayleigh}
\end{equation}
where \(\lambda\) is the relevant wavelength and \(\mathrm{NA}\) is numerical aperture. Equation~\eqref{eq:rayleigh} expresses a criterion for a stated image-formation problem; it is not a universal claim about all estimators, all priors, all sample interactions, or all optical fields. The super-resolution roadmap and a broad fluorescence review both frame modern methods as changes to the conventional measurement model rather than a repeal of wave optics \cite{hell2015,huang2010}.

In particular, the finite aperture still constrains the raw far-field transfer of an ordinary linear image. A super-resolution method earns additional information by controlling what emits, by mixing sample frequencies with a known pattern, by separating emitters in time, by measuring a more informative intensity landscape, by expanding the specimen, or by applying a documented prior. The inference can therefore outperform a conventional separation criterion without asserting that arbitrary unknown structure has appeared from absent data.

\subsection{Four different reported quantities}

The survey finds four quantities that must be reported separately. \emph{PSF width or Rayleigh scale} describes a conventional image-response baseline. \emph{Localization precision} is uncertainty in a fitted emitter coordinate under a declared noise and PSF model. \emph{Structural resolution} asks whether a target separation, width, topology, or dynamic change has been faithfully distinguished. \emph{Coordinate accuracy} includes systematic effects such as drift, registration, refractive-index mismatch, and label linkage. Localization-based microscopy reviews emphasize that precision alone does not establish resolved biological structure \cite{baddeley2018,lelek2021}.

\begin{table*}[t]
\caption{Resolution quantities that must not be substituted for one another.}
\label{tab:quantities}
\centering
\footnotesize
\setlength{\tabcolsep}{3pt}
\begin{tabular}{p{0.20\textwidth} p{0.30\textwidth} p{0.31\textwidth} p{0.13\textwidth}}
\hline
\textbf{Quantity} & \textbf{Question answered} & \textbf{Minimum evidence} & \textbf{Invalid shortcut}\\
\hline
Rayleigh or PSF baseline & What separation is expected from the stated conventional optical model? & Wavelength, NA, contrast criterion, and image-formation assumptions. & Treating it as a universal bound on all measurements.\\
Localization precision & How uncertain is one fitted emitter coordinate? & Photon/background model, PSF model, pixel size, fit diagnostics, and drift handling. & Calling a small coordinate uncertainty structural resolution.\\
Structural resolution & Can the target separation, geometry, or dynamics be distinguished? & Empirical separation recovery, sampling density, label geometry, calibration, and model sensitivity. & Inferring it from a sharp rendering or a nominal instrument specification.\\
Coordinate accuracy & How close is a reported coordinate to the physical target? & Registration, drift, aberration, linkage, and calibration controls. & Reporting precision while omitting systematic offsets.\\
\hline
\end{tabular}
\end{table*}

\section{Measurement-Model Boundary Cases}

\subsection{Band limit, identifiability, and extra information}

A compact statement of conventional optical imaging is \(\mathbf{y}=H\mathbf{x}+\mathbf{n}\), where \(\mathbf{x}\) is the object, \(H\) is the stated measurement operator, \(\mathbf{y}\) is recorded data, and \(\mathbf{n}\) is noise. In a linear far-field image, finite aperture limits the spatial frequencies represented by \(H\). When the object is allowed to be arbitrary and no additional information is supplied, different fine structures can generate indistinguishable data. This is the regime in which a diffraction-limited separation statement has its strongest meaning.

Super-resolution is possible only because the measurement problem is changed or constrained. SIM supplies a known illumination pattern that mixes sample content into a detectable band. STED and RESOLFT control fluorescent states and confine an effective emission response. SMLM adds temporal sparsity and an emitter model. MINFLUX makes an intensity minimum particularly informative for position. Sequential imaging and expansion add sampling rounds or physical scale. These are explicit additions to \(H\), to the accessible data, or to the admissible object class, rather than semantic exceptions to diffraction \cite{hell2015,huang2010,balzarotti2017,reinhardt2023}.

The decisive concept is identifiability. A sub-Rayleigh estimate is scientifically meaningful only when the data and registered assumptions distinguish that parameter from alternatives. A sparse two-emitter model can be identifiable while an arbitrary dense intensity pattern is not. A learned reconstruction can be useful only when its prior is tested on an appropriate distribution; apparent detail is not evidence if another plausible prior produces a different answer from the same raw data. TraceRes consequently requires every claim to expose the object constraints that make its inverse problem solvable and to test those constraints with cases whose relevant truth is known.

\subsection{Boundary cases for a future benchmark}

The benchmark does not ask which method draws the sharpest line. It asks whether the added information channel survives perturbations that remove or weaken its essential assumption. If the SIM pattern is unknown or phase calibration fails, the claimed frequency shift is unsupported. If state control is absent, an effective-PSF claim returns to the ordinary excitation profile. If SMLM density is too high, localization fits cannot establish a dense structural separation. If targeted-localization calibration drifts, photon-efficient precision may not transfer to coordinate accuracy. If expansion is nonuniform, physical rescaling alone does not preserve a molecular distance.

\begin{table*}[t]
\caption{Measurement-model boundary cases that distinguish supported super-resolution from an unsupported extrapolation.}
\label{tab:boundaries}
\centering
\footnotesize
\setlength{\tabcolsep}{3pt}
\begin{tabular}{p{0.19\textwidth} p{0.27\textwidth} p{0.28\textwidth} p{0.18\textwidth}}
\hline
\textbf{Measurement model} & \textbf{Added information or constraint} & \textbf{Required falsification check} & \textbf{Claim boundary}\\
\hline
Conventional linear far field & Finite-NA intensity data only; no controlled state, pattern, sparsity, or validated prior. & Verify wavelength, NA, contrast criterion, and ordinary image model. & Rayleigh or PSF baseline applies to the stated task.\\
SIM or patterned illumination & Known pattern frequency, phase, orientation, modulation, and reconstruction model. & Perturb pattern calibration and test recovery of known separations. & Frequency mixing does not validate an arbitrary reconstruction.\\
STED or RESOLFT & Controlled fluorescent-state transition and measured effective-emission response. & Test depletion profile, background, photophysics, aberration, and sample tolerance. & A smaller effective PSF is not automatically a structure claim.\\
SMLM or targeted localization & Temporal sparsity or intensity-minimum measurement plus emitter/PSF model. & Vary density, photons, blinking, drift, linkage, and fit model against hidden truth. & Precision is limited to the registered coordinate or separation task.\\
Sequential or expansion route & Additional rounds, physical scale factor, registration, and label-retention model. & Test distortion, scale variation, registration, and calibration transfer. & Molecular-scale conclusion is restricted to the validated transformed sample regime.\\
\hline
\end{tabular}
\end{table*}

This analysis gives a direct answer to the title. Diffraction is a genuine limit of an optical transfer problem that is left unchanged. It is not a ban on all sub-Rayleigh inference, because experiments can supply information that an ordinary image does not contain. The scientific burden shifts to showing that the new information is real, calibrated, task-relevant, and sufficient to exclude the alternatives that would otherwise make the fine feature non-identifiable.

\section{Survey Evidence and Usable Gaps}

\subsection{Mechanisms that change the information budget}

Structured illumination microscopy (SIM) uses a known spatial pattern to shift otherwise inaccessible sample frequencies into the detection band. In its common linear form, SIM supplies an approximately two-fold resolution gain relative to conventional wide-field imaging, while nonlinear variants and reconstruction choices require additional care \cite{hell2015,wegel2016}. STED and related RESOLFT methods use controlled fluorescent-state transitions to shrink the effective emitting region. Their practical result depends on photophysics, depletion dose, background, aberration, and sample tolerance rather than on NA alone \cite{hell2015,wegel2016}.

Single-molecule localization microscopy (SMLM), including PALM and STORM families, makes emitting fluorophores sparse in time and estimates the centroids of their diffraction-broadened images. The approach creates an information channel through temporal separation and a PSF model, but it remains limited by photon yield, blinking, density, sampling, reconstruction, drift, and labels \cite{baddeley2018,lelek2021}. Targeted localization methods illustrate another route. MINFLUX experiments use intensity minima to make each detected photon more informative about position; reported nanometer-scale localization remains conditional on fluorophore, stability, optical, and sample settings \cite{balzarotti2017,schmidt2021}.

Sequential-imaging and expansion-enabled approaches further demonstrate that the relevant measurement system includes physical scale, label geometry, and multiple acquisition cycles. Recent reports show molecular-scale or Angstrom-scale information under specific experimental conditions, not an unlimited universal resolving power \cite{reinhardt2023,zwettler2020}. Across modalities, the correct comparison is therefore a task-specific trade-off among spatial scale, temporal scale, depth, field of view, light dose, labeling, robustness, and quantitative validity \cite{valli2021,wegel2016}.

\subsection{Research gaps}

The Survey Agent accepted five linked gaps. First, PSF/Rayleigh width, localization precision, structural resolution, and coordinate accuracy are too often conflated. Second, modality comparisons lack a shared error record that contains photons, sampling, labeling, drift, background, and temporal constraints. Third, quality-control outputs are not always bound to the exact structural claim. Fourth, a reconstruction pipeline needs blinded truth-known cases rather than only favorable images. Fifth, claims of ``breaking diffraction'' must state both the classical constraint bypassed and the constraints that remain. Quality-assessment and reconstruction frameworks make such auditing increasingly feasible but do not eliminate the need for a declared scope \cite{laine2019,gustafsson2016}.

\section{Idea: A Task-Specific Resolution Passport}

\subsection{TraceRes candidate record}

The Idea Agent selected \traceres{} over a physics-only taxonomy and a visual-sharpness score. A taxonomy explains technique families but does not validate a particular image. A sharpness score is rejected because visual contrast is not evidence of structural recovery. \traceres{} instead attaches a structured record to each candidate claim \(c\):
\begin{equation}
 P_c=\{T_c,M_c,I_c,Q_c,E_c,V_c,S_c,F_c\},
\label{eq:passport}
\end{equation}
where \(T_c\) is the declared task, \(M_c\) the modality and acquisition model, \(I_c\) the added information channel, \(Q_c\) calibration and quality data, \(E_c\) the error budget, \(V_c\) blinded validation, \(S_c\) the claim status, and \(F_c\) feasible follow-up actions. The record makes the proposed intervention explicit: it does not modify a microscope; it prevents a conclusion from exceeding what the measurement supports.

\begin{figure*}[t]
\centering
\fbox{\begin{minipage}{0.94\textwidth}
\centering
\small
\textbf{Task and specimen} \(\rightarrow\) \textbf{modality and information channel} \(\rightarrow\) \textbf{calibration, sampling, and error budget} \(\rightarrow\) \textbf{blinded separation test} \(\rightarrow\) \textbf{bounded claim status and follow-up}\\[3pt]
SIM: patterned-frequency mixing. \quad STED/RESOLFT: state-constrained effective PSF. \quad SMLM: temporal sparsity and PSF fitting.\\
MINFLUX: targeted intensity-minimum localization. \quad Sequential/expansion: additional sampling or physical scale.
\end{minipage}}
\caption{TraceRes evidence flow. All branches retain their distinct physics while exposing a shared claim-and-validation contract.}
\label{fig:traceres-flow}
\end{figure*}

\subsection{Bounded status semantics}

The passport may emit five statuses. \emph{Diffraction bounded} applies when the record remains within the conventional imaging model. \emph{Effective PSF confined} applies when nonlinear state control or another physical mechanism supports a narrower effective emission response but structural recovery has not yet been validated. \emph{Localization precision validated} applies to a qualified coordinate uncertainty, not to arbitrary structure. \emph{Structural resolution validated} requires task-matched separation recovery, sampling, calibration, and error controls. \emph{Claim not supported} applies whenever an essential provenance, calibration, sampling, or alternative-explanation field is absent.

These statuses are not cautious synonyms for the same conclusion. They create a useful escalation path. A method can legitimately deliver a highly precise molecular coordinate while lacking sufficient labels to resolve a filament spacing. Conversely, a modest nominal spatial scale may be the strongest scientifically valid choice for a fast, light-sensitive sample if it answers the registered dynamic task with lower artifact risk.

\subsection{Structural claim envelope}

TraceRes proposes a task-specific lower-bound envelope for a released structural scale,
\begin{equation}
 r_{\mathrm{claim}}=\max\{r_{\mathrm{emp}},\;2d_{\mathrm{Nyq}},\;e_{\mathrm{label}},\;2\sigma_{\mathrm{reg}},\;r_{\mathrm{drift}}\},
\label{eq:envelope}
\end{equation}
where \(r_{\mathrm{emp}}\) is an empirical separation or correlation-based estimate, \(d_{\mathrm{Nyq}}\) is sampling spacing, \(e_{\mathrm{label}}\) is label-linkage uncertainty, \(\sigma_{\mathrm{reg}}\) is registration uncertainty, and \(r_{\mathrm{drift}}\) is drift or motion contribution. Equation~\eqref{eq:envelope} is a release gate, not a universal law or replacement for method-specific estimators. It prevents a report from claiming a structural scale smaller than its largest relevant unresolved constraint.

\section{ExperimentDesign: Falsifiable Benchmark Protocol}

\subsection{Hypotheses and design boundary}

H1 tests whether typed passports reduce the rate at which localization precision is reported as structural resolution. H2 tests whether the envelope in \eqref{eq:envelope} improves calibration on blinded truth-known cases compared with nominal-resolution reporting. H3 tests whether source-bound error ledgers reduce omitted artifact and alternative explanations in independent audits. H4 tests whether a passport-constrained modality choice improves expected valid-information yield at a matched resource budget.

The design is strictly `DESIGN ONLY.' It allows planned analysis of authorized public benchmarks, synthetic data, and human-approved physical calibration targets. It does not authorize a microscope, laser, live-cell experiment, fluorophore, biological sample, high-power illumination, reconstruction run, or public performance claim. All actual acquisition and release decisions require domain and institutional review.

\begin{table*}[t]
\caption{Pre-registered TraceRes benchmark phases. No phase has been executed by this proposal.}
\label{tab:design}
\centering
\footnotesize
\setlength{\tabcolsep}{3pt}
\begin{tabular}{p{0.14\textwidth} p{0.31\textwidth} p{0.27\textwidth} p{0.18\textwidth}}
\hline
\textbf{Phase} & \textbf{Planned operation} & \textbf{Primary diagnostic} & \textbf{Permitted conclusion}\\
\hline
0: registration & Freeze task types, truth generators, thresholds, metrics, analysis versions, and release wording. & Manifest completeness and human approval. & No empirical conclusion.\\
1: controlled cases & Generate or obtain authorized truth-known cases spanning separation, density, photons, background, drift, label offset, aberration, and motion. & Hidden truth label and provenance card. & Input suitability only.\\
2: modality records & Bind modality mechanism, acquisition model, PSF/reconstruction model, calibration, and error terms to each case. & Passport field completeness. & Bounded method description only.\\
3: blinded assessment & Evaluate separation recovery and claim status without revealing the expected answer to the primary analyst. & Calibration curve, confusion matrix, and alternative coverage. & Status validation conditional on the registered cases.\\
4: decision audit & Compare single-metric, nominal-resolution, and TraceRes reporting at matched resource assumptions. & False claim rates, valid-information yield, and scope stability. & Prospective method-selection recommendation only.\\
\hline
\end{tabular}
\end{table*}

\subsection{Variables and controls}

Independent variables are modality class, information channel, pair separation, emitter density, photon count, background, drift, label offset, aberration, field of view, and temporal change. Dependent variables are separation recovery, localization precision, coordinate accuracy, structural false-positive and false-negative rates, claim-type confusion, and valid-information yield. Controls include the truth generator, pixel size, wavelength/NA configuration where relevant, raw-data identity, calibration record, analysis version, blinded label, and threshold policy.

The validation set is deliberately counterfactual. It includes a case in which two targets are separated but undersampled, a case in which a model produces a narrow feature from an incorrect prior, a case in which drift mimics displacement, a case in which label linkage dominates coordinate error, and a case in which a fast dynamic sample invalidates an otherwise high-spatial-resolution reconstruction. A method earns a status only if its registered evidence discriminates the corresponding alternative. A favorable image is not a substitute for a passed counterfactual.

\subsection{Task routing and acceptance rules}

No modality is intrinsically the best super-resolution technique. The correct choice begins with the scientific task. A pair-separation task needs a truth-known separation recovery curve. A molecular-coordinate task needs a localization-precision model together with registration and linkage controls. A topology task needs sampling and artifact evidence that distinguish a continuous structure from clustered labels or a reconstruction prior. A dynamics task needs a temporal sampling and light-dose budget that avoids averaging away the event under study. The TraceRes passport routes a proposed measurement to the narrowest claim supported by these requirements rather than to the smallest nominal spatial number.

For a task \(T\), TraceRes defines a prospective valid-information yield
\begin{equation}
 Y(m,T)=\frac{\Pr\!\left(A_T=1\mid m,T,\mathcal{C}\right)\,G_T}{C_m},
\label{eq:yield}
\end{equation}
where \(m\) is a modality, \(A_T\) is the registered task-acceptance event, \(\mathcal{C}\) is the declared calibration and constraint set, \(G_T\) is the information gain relevant to the task, and \(C_m\) is an approved light, time, sample, and analysis cost. Equation~\eqref{eq:yield} is not an empirical ranking result. It defines the comparison that H4 will test. A method with a less extreme spatial number can have higher yield if it is calibrated, sufficiently fast, less perturbative, and better matched to the target.

\begin{table*}[t]
\caption{Task-specific acceptance rules planned for TraceRes.}
\label{tab:routing}
\centering
\footnotesize
\setlength{\tabcolsep}{3pt}
\begin{tabular}{p{0.18\textwidth} p{0.29\textwidth} p{0.30\textwidth} p{0.17\textwidth}}
\hline
\textbf{Scientific task} & \textbf{Minimum acceptance evidence} & \textbf{Counterfactual that must be excluded} & \textbf{Permitted status}\\
\hline
Pair separation or feature width & Empirical recovery across known spacings plus sampling and error envelope. & A single broadened feature, underlabeling, drift, or processing prior mimics two structures. & Structural resolution validated only after passed audit.\\
Molecular coordinate or trajectory & Photon/background fit, localization precision, registration, linkage, and drift controls. & A precise centroid is offset from the physical molecule or changes because of stage/sample motion. & Localization precision validated or coordinate-accuracy scoped.\\
Topology or connectivity & Sampling-density audit, reconstruction stability, and independent feature-preservation test. & Blinking, clustering, sparse labels, or a prior creates apparent branches, holes, or continuity. & Structural status limited to audited topology.\\
Fast live or light-sensitive dynamics & Temporal sampling, photodose, motion, and field-of-view budget tied to event duration. & Averaging, bleaching, or sparse temporal acquisition creates or removes a dynamic transition. & Task-matched dynamic information only.\\
\hline
\end{tabular}
\end{table*}

The routing logic also makes the conventional diffraction baseline productive. If no additional information channel and no validated inverse model are available, a diffraction-bounded result can still be the correct, high-quality outcome. If a mechanism such as SMLM or STED is available but the required task evidence is incomplete, the passport identifies the next calibration rather than encouraging an overstatement. This is a direct experimental distinction: future work can test whether routing by task and acceptance evidence improves the fraction of claims that survive blinded review.

\section{Alternative Explanations and Falsification}

\subsection{Technique-specific counterfactuals}

For SIM, patterned illumination may shift frequencies but a reconstruction can remain sensitive to phase, modulation contrast, out-of-focus light, and algorithmic choices. A TraceRes record requires the pattern model, raw phase series, reconstruction version, and an empirical separation test. For STED or RESOLFT, a nominally narrow effective PSF can be undermined by depletion-beam aberration, background, photobleaching, and sample perturbation; calibration therefore must include the relevant depth and dye regime. For SMLM, sparse centroid fits must be separated from a dense-structure claim by sampling density, blinking, linkage, drift, and reconstruction tests \cite{baddeley2018,lelek2021}.

For MINFLUX, the targeted intensity landscape changes photon efficiency but does not remove the need to verify stability, calibration, labeling, and task transfer \cite{balzarotti2017,schmidt2021}. For expansion and sequential methods, a scale factor or sequential-readout gain must be linked to distortion, label retention, registration, and physical reference structure \cite{reinhardt2023,zwettler2020}. These alternatives are not generic disclaimers. They are the mechanisms by which the proposal can fail, refine its scope, or choose a different modality.

\begin{table*}[t]
\caption{Failure-mode matrix. An ambiguous result narrows a future claim and chooses the next discriminating action.}
\label{tab:failure}
\centering
\footnotesize
\setlength{\tabcolsep}{3pt}
\begin{tabular}{p{0.22\textwidth} p{0.27\textwidth} p{0.26\textwidth} p{0.17\textwidth}}
\hline
\textbf{Detected condition} & \textbf{Why it threatens interpretation} & \textbf{Required TraceRes response} & \textbf{Next action}\\
\hline
Small fitted coordinate uncertainty & Precision does not establish pair separation or structural sampling. & Retain localization-precision status; inspect density and empirical separation. & Add truth-known spacing and Nyquist audit.\\
Sharp reconstruction changes under plausible prior & Model assumptions may create a non-identifiable feature. & Retain claim-not-supported status and publish model sensitivity. & Use blinded alternative-model and raw-data checks.\\
Physical nanoruler passes but sample differs & Calibration may not transfer across depth, label, motion, or scattering regime. & Restrict scope to calibrated regime. & Validate a matched sample class.\\
Drift or registration comparable to claimed scale & Coordinate accuracy is insufficient for relative-structure inference. & Add drift/registration term to envelope. & Acquire approved fiduciary or registration calibration.\\
High spatial scale conflicts with dynamic task & Temporal averaging or photodose can erase the phenomenon of interest. & Compare valid-information yield, not nominal spatial value alone. & Route to a modality with appropriate temporal/light budget.\\
\hline
\end{tabular}
\end{table*}

\section{Interpretive Outcomes and Limits}

If H1--H4 are supported in a future execution, TraceRes would not announce a new absolute diffraction limit. It would establish a more useful practical result: the conditions under which a specific claim is supported, the information resource that enabled it, and the largest remaining constraint. A strong conclusion could state that a task-matched structural separation was recovered under a named modality, calibration object, sampling condition, error envelope, and alternative-model audit. A negative conclusion could state that a given data regime does not support the proposed structural scale, rather than claiming that the structure does not exist.

The limits are substantive. A finite benchmark cannot cover every fluorophore, tissue, material, scattering state, dynamic process, optical aberration, or reconstruction method. Truth-known synthetic data can fail to represent biological complexity. Physical standards can fail to reproduce labeling and motion in living samples. Deep-learning or prior-driven methods require particular scrutiny because the prior can be useful while also becoming a source of non-identifiability. These limits are exactly why the passport records scope and follow-up rather than reporting a context-free number.

The direct answer to the topic is therefore conditional but decisive. There is a diffraction limit for a conventional, stated far-field image-formation model. There is not one universal numerical ceiling on every optical inference. Super-resolution methods can bypass a conventional separation scale by exploiting additional physical control, statistical structure, or sample transformation. Their scientific value is secured not by the phrase ``beyond diffraction,'' but by a task-specific validation that distinguishes true structural information from precision, contrast, or reconstruction appearance.

\section{Conclusion}

Optical super-resolution is most accurately understood as measurement design. The classical diffraction baseline remains indispensable because it states what an ordinary optical image can transfer. SIM, STED/RESOLFT, SMLM, MINFLUX, sequential imaging, and expansion methods add different resources that can make finer information accessible. Each resource also introduces distinct calibration and artifact obligations.

\traceres{} operationalizes this distinction. It ties a scientific claim to a resolution type, information channel, multi-term error budget, blinded validation, bounded status, and follow-up action. This converts an overbroad question about an absolute limit into empirical questions that can be tested, audited, revised, and compared across modalities. The proposed workflow is assertive where evidence is established: conventional diffraction is not the final word on optical inference. It is equally explicit that no asserted nanoscale structure should outrun its sampling, accuracy, and calibration evidence.

\appendices

\section{Minimum Resolution-Passport Schema}

Every future record contains task and feature definition; modality and information channel; raw-data and pipeline versions; wavelength, NA, pixel size, and acquisition conditions where applicable; calibration object; PSF or reconstruction model; photon/background and sampling estimates; label geometry; drift, registration, and aberration controls; declared empirical resolution metric; status; scope; source keys; reviewer decisions; and the next discriminating action. A missing mandatory field prevents escalation to structural-resolution status.

\begin{table*}[t]
\caption{Auditable fields for a minimum TraceRes record. No row is an observed result.}
\label{tab:passport-schema}
\centering
\footnotesize
\setlength{\tabcolsep}{3pt}
\begin{tabular}{p{0.21\textwidth} p{0.34\textwidth} p{0.33\textwidth}}
\hline
\textbf{Field block} & \textbf{Required content} & \textbf{Audit and release rule}\\
\hline
Task and information channel & Feature type, desired separation or dynamics, modality, optical or statistical resource beyond the baseline model. & A claim cannot be evaluated if its target quantity or added information channel is unspecified.\\
Provenance and calibration & Data and pipeline identity, acquisition model, nanoruler or truth generator, parameter registry, and reviewer identity. & Version ambiguity or uncalibrated change forces claim-not-supported status.\\
Error and sampling ledger & Photon/background model, empirical separation metric, label density/linkage, drift, registration, aberration, motion, and temporal averaging. & The released structural scale cannot be smaller than the largest relevant unexcluded term.\\
Validation and decision card & Blinded-case outcome, model sensitivity, remaining alternatives, status, allowed wording, and follow-up value. & Status changes require a new versioned card and independent human review.\\
\hline
\end{tabular}
\end{table*}

\section{Evidence-to-Claim Traceability}

The final claim is intentionally narrower than a generic statement that diffraction has been defeated. The Survey sources support the conditional conventional baseline and the mechanism taxonomy; the Idea Agent contributes the passport as a proposed integration; the ExperimentDesign Agent specifies how a future benchmark can try to falsify the proposed integration. TraceRes therefore treats an evidence source, a calibration object, and a final scientific conclusion as different objects. A source may justify why a mechanism is plausible, while only a task-matched future validation can justify a new structural claim for a particular sample.

This separation protects both scientific ambition and reproducibility. It permits a direct claim that targeted localization, nonlinear state control, patterned illumination, temporal sparsity, and expansion can provide information inaccessible to an ordinary diffraction-bounded image. It prohibits the stronger but unsupported claim that any one of these mechanisms guarantees a stated biological structure. For every future publication, the traceability record must name the exact modality, information channel, validation scope, maximum remaining error term, and allowed wording. An independent reviewer can then test the claim without reconstructing the entire research history from prose.

\begin{table*}[t]
\caption{Evidence-to-claim traceability matrix for the proposed report and future TraceRes releases.}
\label{tab:traceability}
\centering
\footnotesize
\setlength{\tabcolsep}{3pt}
\begin{tabular}{p{0.20\textwidth} p{0.25\textwidth} p{0.27\textwidth} p{0.18\textwidth}}
\hline
\textbf{Claim class} & \textbf{Registered evidence basis} & \textbf{Required future validation} & \textbf{Prohibited wording}\\
\hline
Conventional diffraction baseline & Stated far-field model, NA, wavelength, contrast criterion, and the mechanism reviews \cite{hell2015,huang2010}. & Verify that the experiment actually remains within the declared conventional model. & ``Diffraction prevents every sub-Rayleigh inference.''\\
Mechanism-enabled information gain & Method-specific sources for SIM, STED/RESOLFT, SMLM, MINFLUX, sequential, or expansion routes \cite{wegel2016,lelek2021,schmidt2021,zwettler2020}. & Bind the actual information channel and configuration to the passport. & ``A sharp image proves structure regardless of method.''\\
Localization or coordinate claim & Photon/background model, fit diagnostics, registration, drift, and linkage record \cite{baddeley2018,lelek2021}. & Recover known coordinates and offsets in a matched calibration regime. & ``Localization precision equals structural resolution.''\\
Structural or topology claim & Sampling, empirical separation test, reconstruction sensitivity, calibration, and the envelope in \eqref{eq:envelope}. & Pass blinded truth-known cases and disclose residual constraints. & ``The sample is resolved because the nominal instrument value is small.''\\
Technique-selection recommendation & Task objective, light/time/sample cost, counterfactuals, and prospective yield in \eqref{eq:yield}. & Compare registered alternatives at matched approved resources. & ``The numerically finest method is universally best.''\\
\hline
\end{tabular}
\end{table*}

A traceability matrix also prevents the Author stage from inventing empirical content. The author can synthesize why the status categories are needed and how a future test would discriminate alternatives, but cannot claim that the test has already passed. Under this contract, a negative or ambiguous result is not a failed study: it is a bounded result that identifies the calibration, sampling, model, or physical control required for the next valid escalation.

\section{Proposal Status and Source Boundary}

This document is a research-plan artifact derived from the Survey, Idea, and ExperimentDesign handoffs in the same project directory. The sources are registered in the Survey manifest. Proposed contributions are the TraceRes passport, hypotheses, counterfactual validation protocol, decision rule, and release schema. No microscope acquisition, benchmark run, biological observation, or super-resolution performance result is claimed. The requested in-app browser could not connect because of a Windows ACL failure; source metadata were cross-validated through OpenAlex and AnySearch, while direct publisher-page verification remains a human-review task.

\begin{thebibliography}{00}
\bibitem{hell2015} S. W. Hell et al., ``The 2015 super-resolution microscopy roadmap,'' \emph{Journal of Physics D: Applied Physics}, vol. 48, 443001, 2015, doi: 10.1088/0022-3727/48/44/443001.
\bibitem{huang2010} B. Huang, H. P. Babcock, and X. Zhuang, ``Breaking the diffraction barrier: Super-resolution imaging of cells,'' \emph{Cell}, vol. 143, no. 7, pp. 1047--1058, 2010, doi: 10.1016/j.cell.2010.12.002.
\bibitem{baddeley2018} D. Baddeley and J. Bewersdorf, ``Biological insight from super-resolution microscopy: What we can learn from localization-based images,'' \emph{Annual Review of Biochemistry}, vol. 87, pp. 965--989, 2018, doi: 10.1146/annurev-biochem-060815-014801.
\bibitem{valli2021} J. Valli et al., ``Seeing beyond the limit: A guide to choosing the right super-resolution microscopy technique,'' \emph{Journal of Biological Chemistry}, vol. 297, 100791, 2021, doi: 10.1016/j.jbc.2021.100791.
\bibitem{wegel2016} E. Wegel et al., ``Imaging cellular structures in super-resolution with SIM, STED and localisation microscopy: A practical comparison,'' \emph{Scientific Reports}, vol. 6, 27290, 2016, doi: 10.1038/srep27290.
\bibitem{lelek2021} M. Lelek et al., ``Single-molecule localization microscopy,'' \emph{Nature Reviews Methods Primers}, vol. 1, 39, 2021, doi: 10.1038/s43586-021-00038-x.
\bibitem{balzarotti2017} F. Balzarotti et al., ``Nanometer resolution imaging and tracking of fluorescent molecules with minimal photon fluxes,'' \emph{Science}, vol. 355, no. 6325, pp. 606--612, 2017, doi: 10.1126/science.aak9913.
\bibitem{schmidt2021} R. Schmidt et al., ``MINFLUX nanometer-scale 3D imaging and microsecond-range tracking on a common fluorescence microscope,'' \emph{Nature Communications}, vol. 12, 1478, 2021, doi: 10.1038/s41467-021-21652-z.
\bibitem{reinhardt2023} S. Reinhardt et al., ``\AA ngstrom-resolution fluorescence microscopy,'' \emph{Nature}, vol. 617, pp. 711--716, 2023, doi: 10.1038/s41586-023-05925-9.
\bibitem{zwettler2020} F. U. Zwettler et al., ``Molecular resolution imaging by post-labeling expansion single-molecule localization microscopy,'' \emph{Nature Communications}, vol. 11, 3388, 2020, doi: 10.1038/s41467-020-17086-8.
\bibitem{laine2019} R. F. Laine et al., ``NanoJ: A high-performance open-source super-resolution microscopy toolbox,'' \emph{Journal of Physics D: Applied Physics}, vol. 52, 163001, 2019, doi: 10.1088/1361-6463/ab0261.
\bibitem{gustafsson2016} N. Gustafsson et al., ``Fast live-cell conventional fluorophore nanoscopy with ImageJ through super-resolution radial fluctuations,'' \emph{Nature Communications}, vol. 7, 12471, 2016, doi: 10.1038/ncomms12471.
\end{thebibliography}

\end{document}
